\chapter{Methoden}
\label{chap:methoden}

\section{Hardware-Aufbau}

\subsection{Raspberry Pi 5 und Peripherie}

Der Raspberry Pi\,5 (8\,GB RAM) dient als Rechenplattform. Der Hailo-8-NPU
(AI HAT+, 27\,TOPS) wird über PCIe angebunden.

\subsection{Pi HQ Camera (IMX477) mit angepasster Optik}

Die Pi HQ Camera nutzt einen C-Mount und einen angepassten Tubus mit einem
125--150\,mm Achromat (ThorLabs AC127-125-A bzw.\ AC127-150-A) anstelle des
Standard-50\,mm-Objektivs. Dadurch wird die Sensornutzung von ~25\,\% auf
~85--95\,\% erhöht (siehe Abschnitt Optik-Kalkulation).

\subsection{OpenFlexure-Mikroskop}

Das 3D-gedruckte Gehäuse wird mit der \texttt{hq\_camera}-Variante von
OpenFlexure\,v7 erstellt. Die Montage folgt den offiziellen Assembly-Anleitungen.

\section{Software-Infrastruktur}

\subsection{Ansible-Automatisierung}

Die gesamte Software wird über Ansible-Playbooks bereitgestellt: Basissystem,
Hailo-Treiber, OpenFlexure Server\,v3 und Motorsteuerung.

\subsection{OpenFlexure Server v3}

Der Server (FastAPI/LabThings) steuert Kamera und (optional) Motorstage über
die Weboberfläche auf Port\,5000.

\section{KI-Pipeline}

\subsection{Datensammlung}

Zeitraffer-Aufnahmen über die OpenFlexure-API oder direkt via picamera2.

\subsection{Annotation}

CVAT-Server auf separater VPS für semantische Segmentierung von Osteoblasten,
epithelialen Zellen und Kalziumablagerungen.

\subsection{Training}

YOLOv8s-seg (Ultralytics) mit angereicherten Daten, exportiert als ONNX und
kompiliert zu HEF für den Hailo-8-NPU.

\subsection{Inferenz und Tracking}

Echtzeit-Inferenz auf Hailo-8, Zellverfolgung mit ByteTrack, morphologische
Analyse (Fläche, Perimeter, Zirkularität, Migrationsgeschwindigkeit).